\documentclass[11pt]{article}
\usepackage[a4paper,margin=27mm]{geometry}
\usepackage{amsmath,amssymb,amsthm}
\usepackage{longtable,array}
\usepackage[hidelinks]{hyperref}
\hypersetup{pdftitle={Classification of order-ten binary 3-frameproof codes},pdfauthor={GPT-6 Astra}}
\urlstyle{tt}
\setlength{\emergencystretch}{2em}
\newtheorem{theorem}{Theorem}[section]
\newtheorem{lemma}[theorem]{Lemma}
\newtheorem{proposition}[theorem]{Proposition}
\newtheorem{corollary}[theorem]{Corollary}
\theoremstyle{definition}
\newtheorem{definition}[theorem]{Definition}
\theoremstyle{remark}
\newtheorem{remark}[theorem]{Remark}
\newcommand{\bits}{\{0,1\}}
\newcommand{\I}{I}
\newcommand{\shf}{\operatorname{SHF}}
\newcommand{\wt}{\operatorname{wt}}
\newcommand{\lexle}{\leq_{\mathrm{lex}}}
\newcommand{\lexlt}{<_{\mathrm{lex}}}
\newcommand{\ind}[1]{\mathbf{1}_{\{#1\}}}
\title{Classification of order-ten binary 3-frameproof codes\\
\large A computer-assisted proof formalized in Lean}
\author{GPT-6 Astra}
\date{8 October 2026}

\begin{document}
\maketitle


\begin{abstract}
Every binary $\shf(10;10,2,\{1,3\})$ is a permutation matrix up to
independent complementation of rows. This resolves the order-ten case
explicitly left open by Guo, Stinson and van Trung in 2015. The proof
reduces a hypothetical nontrivial example to one of four canonical
Boolean formulas on the entries of a $10\times10$ matrix. The reduction
uses an elementary structural lemma, row complementation, and minimization
over a finite permutation orbit. Each formula has 840 separation
conditions, 80 row-order conditions, 63 column-order conditions, and 19
fixed-entry conditions. Lean verifies the analytic reduction, the semantic
Boolean bridges, all four finite refutations, and their integration into
the classification theorem. The finite refutations replay saved LRAT
certificates through the standard \texttt{bv\_check} procedure; their trust
boundary includes Lean's compiled evaluation of the verified checker.
We give the exact formulas, the mathematical proofs behind their
normalization, and instructions for reproducing the formal result.
\end{abstract}

\section{The question and its precise scope}
\label{sec:scope}

Write $\I=\{0,1,\ldots,9\}$ and identify the Boolean values false and true
with $0$ and $1$. A binary matrix is a function
$A:\I\times\I\to\bits$. Rows are coordinates, or hash functions, and
columns are codewords, or elements of the hash-family domain.

\begin{definition}
\label{def:separation}
A row $r$ \emph{separates} a column $a$ from columns $b,c,d$ if
\begin{equation}
\label{eq:separation}
 A_{ra}\ne A_{rb},\qquad A_{ra}\ne A_{rc},\qquad A_{ra}\ne A_{rd}.
\end{equation}
The matrix $A$ has the \emph{separation property} if, for every four
pairwise distinct columns $a,b,c,d\in\I$, some row separates $a$ from
$b,c,d$.
\end{definition}

This is the matrix definition of $\shf(10;10,2,\{1,3\})$ used in
\cite{designs2015}. In a binary row, the three entries on the right of
\eqref{eq:separation} necessarily agree, although the definition does not
need to state their agreement separately. Equivalently, no column can be
reproduced coordinate by coordinate by choosing at each coordinate an
entry from any fixed three other columns. The version of the frameproof
definition allowing coalitions of size at most three is equivalent here:
a smaller coalition can be extended to three other columns, and its
coordinatewise descendants are contained in those of the larger
coalition. There are enough columns to make this extension. Distinctness
of columns is also forced by separation, since any two equal columns
would prevent the required separation of one of them from a coalition
containing the other.

\begin{definition}
\label{def:trivial}
The matrix $A$ is \emph{trivial} if there are a permutation $p$ of $\I$
and bits $\epsilon_r\in\bits$, one for each row, such that
\begin{equation}
\label{eq:trivial}
 A_{rc}=\ind{p(r)=c}\mathbin{\oplus}\epsilon_r
 \qquad(r,c\in\I),
\end{equation}
where $\oplus$ denotes exclusive or.
\end{definition}

Thus trivial means a permutation matrix up to independent row
complementation. Additional row and column permutations are already
absorbed in $p$. Conversely, every trivial matrix has the separation
property: for designated column $a$, choose the unique row $r$ with
$p(r)=a$. Its entry in column $a$ differs from every other entry in that
row, and complementation preserves inequalities.

\begin{theorem}[Classification]
\label{thm:main}
Every binary $10\times10$ matrix with the separation property is trivial.
Equivalently, no nontrivial $\shf(10;10,2,\{1,3\})$ exists.
\end{theorem}

The conclusion is an equality classification, not merely a bound on the
number of codewords. The 2015 bound of Guo, Stinson and van Trung gives
the corresponding uniqueness statement for the relevant orders
$4\leq v\leq9$
\cite{bounds2015}. Their paper on symmetric designs exhibits a
nontrivial order-eleven example from a symmetric $(11,5,2)$ design and
explicitly leaves order ten open in its concluding problem 3
\cite{designs2015}. Neither the problem nor the classification
assumes that the matrix comes from a design, has constant row or column
weight, or has a prescribed automorphism group.

The later binary wide-sense frameproof-code results of Zhao and Zhang
\cite{wide2024} do not by themselves settle this equality case. Their
quoted range $n\leq3t$ reaches nine when $t=3$, while the range
$n<((15+\sqrt{33})/24)(t-1)^2$ does not reach ten. This is an
applicability check on those results, not a claim that they establish the
continued openness of the exact question. A literature audit completed
on 8 October 2026 did not locate an earlier resolution. Its coverage is
incomplete, including some inaccessible full texts; it cannot guarantee
priority or exclude an unpublished result. In particular, the full text
of \cite{sunwang2025} was not available in that audit.

\section{An elementary structural lemma}
\label{sec:structure}

\begin{definition}
A row $r$ has a \emph{private bit at column $a$} if
$A_{rb}\ne A_{ra}$ for every $b\ne a$. A row is \emph{elementary} if
it is constant or has a private bit.
\end{definition}

The binary alphabet is essential in the following argument: all entries
different from a given bit are equal to one another.

\begin{lemma}
\label{lem:private-unique}
A row has a private bit at at most one column.
\end{lemma}
\begin{proof}
Suppose a row has private bits at distinct columns $a$ and $b$. Choose
$c\notin\{a,b\}$, which is possible because there are ten columns.
Privacy at $a$ implies that both $A_{rb}$ and $A_{rc}$ are the unique
binary value different from $A_{ra}$, so $A_{rb}=A_{rc}$. This
contradicts privacy at $b$.
\end{proof}

\begin{lemma}
\label{lem:elementary-separator}
Suppose $b\ne c$ and an elementary row $r$ satisfies
$A_{ra}\ne A_{rb}$ and $A_{ra}\ne A_{rc}$. Then $r$ has its private
bit at $a$.
\end{lemma}
\begin{proof}
The row cannot be constant. Let $e$ be a column at which its bit is
private. If $e=a$, the conclusion holds. Otherwise all entries outside
$e$, including the entry at $a$, have the common bit different from
$A_{re}$. The inequality at $b$ forces $b=e$, and the inequality at
$c$ forces $c=e$. This contradicts $b\ne c$.
\end{proof}

\begin{theorem}[Elementary-row classification]
\label{thm:elementary}
If $A$ has the separation property and every row of $A$ is elementary,
then $A$ is trivial.
\end{theorem}
\begin{proof}
For each column $a$, choose three distinct other columns $b,c,d$.
Separation supplies a row $f(a)$ separating $a$ from those columns.
Lemma~\ref{lem:elementary-separator}, applied to $b,c$, shows that
$f(a)$ has a private bit at $a$.

The map $f:\I\to\I$ is injective: if $f(a)=f(a')$, that row has
private bits at $a$ and $a'$, so
Lemma~\ref{lem:private-unique} gives $a=a'$. An injection between these
two finite sets of size ten is a bijection. Put $p=f^{-1}$. Every row
$r$ therefore has a private bit at column $p(r)$; in particular, the
argument has also ruled out unused constant rows.

Define $\epsilon_r=1\oplus A_{r,p(r)}$. At $c=p(r)$,
$1\oplus\epsilon_r=A_{r,p(r)}$. At every $c\ne p(r)$, privacy and
the binary alphabet give $A_{rc}=1\oplus A_{r,p(r)}=\epsilon_r$.
These two cases are precisely \eqref{eq:trivial}.
\end{proof}

\begin{corollary}
\label{cor:nontrivial-row}
A nontrivial matrix with the separation property has a row containing
at least two zeros and at least two ones.
\end{corollary}
\begin{proof}
By the contrapositive of Theorem~\ref{thm:elementary}, some row is
neither constant nor private at any column. Both bits occur because it
is nonconstant. Neither occurs exactly once, because that occurrence
would be a private bit. Thus both occur at least twice.
\end{proof}

No bound from the earlier classification through order nine is used in
this structural proof.

\section{Normalization and anchoring}
\label{sec:normalization}

\begin{lemma}[Invariance of separation]
\label{lem:invariance}
The separation property is preserved by independent complementation of
rows and by arbitrary permutations of rows and columns.
\end{lemma}
\begin{proof}
For bits $x,y,e$, one has $x\oplus e\ne y\oplus e$ if and only if
$x\ne y$. Consequently, complementing a row preserves each inequality
in \eqref{eq:separation}. A row permutation moves a witness to its
inverse image. A column permutation sends a quartet of distinct
columns to a quartet of distinct columns, to which the original
separation property applies. Applying the inverse row permutation to
the resulting witness proves the assertion.
\end{proof}

For a row $r$, put $C_r=\{c:A_{rc}=1\}$ and
$\wt(A_r)=|C_r|$. Complementing the row replaces $C_r$ by its
complement in $\I$. For definiteness, use the following orientation:
retain $C_r$ if $|C_r|<5$, or if $|C_r|=5$ and $9\notin C_r$;
otherwise replace $C_r$ by its complement. Denote the oriented set by
$\widehat C_r$. Exactly one member of a complementary pair satisfies
this rule, including when both members have size five. In particular,
$|\widehat C_r|\leq5$ for every row.

\begin{lemma}[A distinguished first row]
\label{lem:first-row}
If $A$ is a nontrivial matrix with the separation property, then after
row complements and row and column permutations it has a representative
$N$ and an integer $k\in\{2,3,4,5\}$ such that
\begin{equation}
\label{eq:first-row}
 N_{0c}=\ind{c<k}\qquad(c\in\I).
\end{equation}
One may additionally arrange that every row of $N$ has weight at most
five and has weight at most one or at least $k$.
\end{lemma}
\begin{proof}
Orient every row as above. By Corollary~\ref{cor:nontrivial-row}, at
least one oriented row has weight between two and five: complementation
preserves the assertion that both bits occur at least twice. Choose an
oriented row of smallest weight among those of weight at least two, and
call its weight $k$. This gives $2\leq k\leq5$ and the claimed
restrictions on the weights of all other rows.

Move the chosen row to index zero. Let $C$ be its set of ones, so
$|C|=k$. Choose any bijection from $\{0,\ldots,k-1\}$ to $C$, and
any bijection from $\{k,\ldots,9\}$ to $\I\setminus C$. Their union
is a permutation $p$ of $\I$. Replacing column $c$ by column $p(c)$
gives \eqref{eq:first-row}. Column permutations preserve row weights,
and Lemma~\ref{lem:invariance} preserves separation throughout.
\end{proof}

\begin{lemma}[Anchored representative]
\label{lem:anchor}
The representative in Lemma~\ref{lem:first-row} can be changed by
further row complements to a matrix $B$ with the separation property,
the same first row \eqref{eq:first-row}, and
\begin{equation}
\label{eq:anchor}
 B_{r0}=0\qquad(1\leq r\leq9).
\end{equation}
\end{lemma}
\begin{proof}
Leave row zero unchanged. For $r>0$, replace row $r$ by
$B_{rc}=N_{rc}\oplus N_{r0}$. Its entry at column zero is
$N_{r0}\oplus N_{r0}=0$. Separation is preserved by
Lemma~\ref{lem:invariance}, and the first row is unchanged.
\end{proof}

\begin{remark}
\label{rem:no-weight}
After this anchoring step, no bound on the weights of rows $1,\ldots,9$
is retained or used. Complementing a light row can make it heavy. The
only weight specification needed subsequently is the explicit first
row in \eqref{eq:first-row}. In particular, the finite formulas below
contain no additional row-weight or minimum-weight conditions.
\end{remark}

\section{A simultaneous lexicographic representative}
\label{sec:lex}

Use the order $0<1$. For two finite binary sequences $u,v$ of the same
length, write $u\lexle v$ for the usual lexicographic order, comparing
coordinates in increasing index order. Equivalently,
\begin{equation}
\label{eq:prefix-order}
 \text{for each index }c,\quad
 \bigl(\forall t<c,\ u_t=v_t\bigr)\ \Longrightarrow\ u_c\leq v_c.
\end{equation}
Indeed, if the sequences are distinct, the condition at their first
unequal coordinate says that its pair of bits is $(0,1)$. Conversely,
the first unequal coordinate of a lexicographically ordered pair
prevents a later coordinate from having an equal prefix.

Fix $k\in\{2,3,4,5\}$. The residual row permutations fix row zero.
The residual column permutations fix column zero and permute separately
the two sets
\begin{equation}
\label{eq:blocks}
 U_k=\{1,\ldots,k-1\},\qquad V_k=\{k,\ldots,9\}.
\end{equation}
Both sets are nonempty. These permutations preserve
\eqref{eq:first-row} and \eqref{eq:anchor}.

\begin{lemma}[Finite-orbit double lexicographic ordering]
\label{lem:double-lex}
An anchored matrix $B$ satisfying \eqref{eq:first-row} has a
representative $D$ under the residual permutations such that, writing
$D_c^-=(D_{1c},\ldots,D_{9c})$ for a truncated column,
\begin{align}
\label{eq:row-order}
 D_i&\lexle D_j &&(1\leq i<j\leq9),\\
\label{eq:column-order}
 D_c^-&\lexle D_d^-
 &&(c<d,\ c,d\in U_k\text{ or }c,d\in V_k).
\end{align}
The first row and the anchor remain unchanged.
\end{lemma}
\begin{proof}
The orbit under the residual permutations is finite and nonempty.
Flatten each matrix in row-major order,
\[
 (D_{00},D_{01},\ldots,D_{09},D_{10},\ldots,D_{99}),
\]
and choose a lexicographically least flattened matrix $D$ in the orbit.

If \eqref{eq:row-order} failed for $i<j$, then $D_j\lexlt D_i$.
Swap rows $i$ and $j$, an allowed residual permutation. All rows before
$i$ remain unchanged. Within row $i$, the first difference from the
original matrix changes a $1$ to a $0$, and every preceding entry is
unchanged. The new flattened matrix is smaller, a contradiction.

Suppose \eqref{eq:column-order} failed for a permitted pair $c<d$.
Let $r>0$ be the first row at which these truncated columns differ.
Failure means $D_{rc}=1$ and $D_{rd}=0$, and
$D_{sc}=D_{sd}$ for $1\leq s<r$. Also $D_{0c}=D_{0d}$ because the
two columns belong to the same first-row block. Swap columns $c,d$,
which is an allowed permutation. Every entry in rows before $r$ is
unchanged: the swapped entries were equal. In row $r$, entries before
column $c$ are unchanged, while the entry at $c$ changes from $1$ to
$0$. This again strictly decreases the flattened matrix, a
contradiction. Finally, all orbit elements satisfy the fixed first row
and zero anchor by the definition of the residual group.
\end{proof}

This argument establishes simultaneous ordering in one finite orbit;
it does not assume that separately sorting rows and columns leaves
the other ordering intact.

\begin{theorem}[Canonical reduction]
\label{thm:canonical}
If a nontrivial matrix with the separation property exists, then for
some $k\in\{2,3,4,5\}$ there exists a matrix $D$ with the separation
property satisfying \eqref{eq:first-row}, \eqref{eq:anchor},
\eqref{eq:row-order}, and \eqref{eq:column-order}.
\end{theorem}
\begin{proof}
Apply Lemma~\ref{lem:first-row}, then Lemma~\ref{lem:anchor}, and
then Lemma~\ref{lem:double-lex}. Lemma~\ref{lem:invariance}
preserves the separation property at each step. The last step fixes
the first row and anchor. There are no other conditions on $D$.
\end{proof}

\section{The four explicit Boolean formulas}
\label{sec:formulas}

For each fixed $k\in\{2,3,4,5\}$, introduce the 100 Boolean variables
$x_{rc}$ with $r,c\in\I$. Every empty conjunction below has value
true. An equality between Boolean expressions is interpreted as its
Boolean equality test when it occurs inside a formula.

\subsection{Separation: 840 conditions}
For $a\in\I$ and $b<c<d$ in $\I\setminus\{a\}$ define
\begin{equation}
\label{eq:sep-formula}
 S_{a;bcd}(x)=
 \bigvee_{r=0}^{9}
 \bigl((x_{ra}\oplus x_{rb})\wedge
       (x_{ra}\oplus x_{rc})\wedge
       (x_{ra}\oplus x_{rd})\bigr).
\end{equation}
There are $10\binom{9}{3}=840$ such conditions, each required to be
true. Ordering $b,c,d$ removes duplicates only: separation is symmetric
in these three opposing columns. The designated column $a$ is not
interchanged with them or omitted from the indexing.

\subsection{The first row and the anchor: 19 conditions}
Require the ten equations and nine equations
\begin{equation}
\label{eq:fixed-formula}
 x_{0c}=\ind{c<k}\quad(0\leq c\leq9),\qquad
 x_{r0}=0\quad(1\leq r\leq9).
\end{equation}
In particular $x_{00}=1$, so the anchor does not include row zero.

\subsection{Adjacent row ordering: 80 conditions}
For $1\leq i\leq8$ and $0\leq c\leq9$, put
\begin{align}
 P_{i,c}(x)&=\bigwedge_{0\leq t<c}(x_{it}=x_{i+1,t}),\nonumber\\
\label{eq:row-formula}
 R_{i,c}(x)&=\neg P_{i,c}(x)\ \vee\ \neg x_{ic}\ \vee\ x_{i+1,c}.
\end{align}
Require all $8\cdot10=80$ formulas $R_{i,c}$ to be true. By
\eqref{eq:prefix-order}, these say that each adjacent pair of free
rows is lexicographically ordered. Transitivity then orders all free
rows. The stronger all-pairs conclusion of
Lemma~\ref{lem:double-lex} directly supplies every required adjacent
condition.

\subsection{Adjacent residual column ordering: 63 conditions}
The permitted left endpoints of adjacent column pairs are
\begin{equation}
\label{eq:J}
 J_k=\{1,\ldots,k-2\}\ \cup\ \{k,\ldots,8\},
\end{equation}
where the first set is empty for $k=2$. These are precisely the pairs
$(c,c+1)$ lying inside $U_k$ or $V_k$. They neither compare column zero
with another column nor compare across the block boundary. For
$c\in J_k$ and $1\leq r\leq9$ define
\begin{align}
 Q_{c,r}(x)&=\bigwedge_{1\leq s<r}(x_{sc}=x_{s,c+1}),\nonumber\\
\label{eq:column-formula}
 L_{c,r}(x)&=\neg Q_{c,r}(x)\ \vee\ \neg x_{rc}\ \vee\ x_{r,c+1}.
\end{align}
Require every $L_{c,r}$ to be true. Since
$|J_k|=(k-2)+(9-k)=7$, there are exactly $7\cdot9=63$
conditions in every case. These compare the truncated columns with
row zero omitted, in agreement with \eqref{eq:column-order}.

\subsection{The combined formula and its exact meaning}
Let $F_k(x)$ be the conjunction of the 19 equations in
\eqref{eq:fixed-formula}. Define
\begin{equation}
\label{eq:Phi}
 \begin{aligned}
 \Phi_k(x)={}&F_k(x)\ \wedge\
 \bigwedge_{\substack{a\in\I\\ b<c<d,\ b,c,d\ne a}}S_{a;bcd}(x)\\
 &{}\wedge\ \bigwedge_{\substack{1\leq i\leq8\\0\leq c\leq9}}R_{i,c}(x)
 \ \wedge\ \bigwedge_{\substack{c\in J_k\\1\leq r\leq9}}L_{c,r}(x).
 \end{aligned}
\end{equation}
There are $840+80+63+19=1002$ high-level conjuncts. This is a count
before constant simplification. Some become tautologies after the
fixed entries are substituted. It is not a count of clauses in a
particular auxiliary-variable CNF encoding: for example,
\eqref{eq:sep-formula} contains conjunctions inside a disjunction.

\begin{proposition}[Exact reduction of existence]
\label{prop:exact}
A nontrivial binary $\shf(10;10,2,\{1,3\})$ exists if and only if at
least one of $\Phi_2,\Phi_3,\Phi_4,\Phi_5$ is satisfiable.
\end{proposition}
\begin{proof}
For the forward implication, apply Theorem~\ref{thm:canonical} and
assign $x_{rc}=D_{rc}$. Its separating witnesses make every
$S_{a;bcd}$ true. The first row and anchor give $F_k$, and the
lexicographic prefix implications give $R_{i,c}$ and $L_{c,r}$.
Thus $\Phi_k$ is true.

Conversely, interpret a satisfying assignment as a binary matrix $X$.
For any designated column $a$ and any three distinct other columns,
sort those three indices as $b<c<d$. The corresponding
$S_{a;bcd}$ gives a row separating $a$ from all three, independently
of their order. Therefore $X$ has the separation property. Its first
row has exactly $k$ ones by $F_k$. In a trivial matrix every row has
weight either one or nine, as follows immediately from
\eqref{eq:trivial}. Since $2\leq k\leq5$, $X$ is nontrivial.
\end{proof}

In particular, the symmetry conditions have not silently removed every
representative of a possible counterexample. The forward implication
is established by a constructive sequence of transformations followed
by a minimum over a finite orbit.

\section{The finite refutations and the classification theorem}
\label{sec:finite-check}

\begin{lemma}[Finite Boolean refutations]
\label{lem:finite-unsat}
For each $k\in\{2,3,4,5\}$, the formula $\Phi_k$ is unsatisfiable:
\begin{equation}
\label{eq:finite-unsat}
 \forall x\in\bits^{\I\times\I},\qquad\neg\Phi_k(x).
\end{equation}
\end{lemma}
\begin{proof}[Computer-assisted proof]
The Lean development proves the semantic interpretation of every
conjunct in \eqref{eq:Phi}. It expands the resulting finite Boolean
expression and canonicalizes the concrete matrix indices using
checked definitional equality. The four expressions are refuted by
replaying the saved certificates
\path{certificates/case2.lrat},
\path{certificates/case3.lrat},
\path{certificates/case4.lrat}, and
\path{certificates/case5.lrat} with Lean's standard \path{bv_check}
procedure. Each replay proves the corresponding Boolean contradiction;
no solver answer is inserted as an unproved mathematical hypothesis.
Appendix~\ref{app:computation} specifies the expression construction,
normalization, and certificate method. Section~\ref{sec:trust} states
its compiled-evaluation trust boundary and reproduction commands.
\end{proof}

\begin{proof}[Proof of Theorem~\ref{thm:main}]
Suppose $A$ has the separation property but is not trivial.
Proposition~\ref{prop:exact} supplies $k\in\{2,3,4,5\}$ and a
satisfying assignment of $\Phi_k$. This contradicts
Lemma~\ref{lem:finite-unsat} for that value of $k$. Hence $A$ is trivial.
\end{proof}

\subsection{Exact interface used by the Lean case proofs}
The development uses a direct matrix interface. For each
$k\in\{2,3,4,5\}$, the theorem
\path{anchored_case_k} takes the following arguments: a binary matrix
$A$, its separation property, its fixed first row, its zero anchor,
the all-pairs ordering of free rows, and the ordering of residual
columns. Its conclusion is false. Each such theorem obtains the
1002 concrete Boolean facts through proved semantic lemmas and then
refutes their conjunction.

The Boolean row-prefix implementation is the explicit finite formula
\[
 \bigwedge_{t=0}^{9}\bigl((c\leq t)\vee(x_t=y_t)\bigr).
\]
For $t\geq c$ the first disjunct is true; for $t<c$ it reduces to
$x_t=y_t$. A checked lemma identifies this formula with equality on
all earlier coordinates. The column-prefix version uses $t=1,\ldots,9$
and likewise omits row zero. The Boolean separation implementation
expands the existential row witness into exactly ten disjuncts. Thus
these helper definitions have the meanings of
\eqref{eq:sep-formula}, \eqref{eq:row-formula}, and
\eqref{eq:column-formula}, not merely similar-looking encodings.

The formal canonicalization theorem produces an integer $k$ satisfying
$2\leq k\leq5$ together with precisely the five matrix properties
used by the case theorems. Case analysis on this interval gives the
four alternatives. The final integration proves the unconditional statements
\begin{verbatim}
every_order_ten_shf_is_trivial
  (A : BinaryMatrix 10) (h : IsSHF A) : IsTrivial A

not_openProblem10 : Not OpenProblem10
\end{verbatim}
Here \path{OpenProblem10} is the definition
$\exists A,\ \mathrm{IsSHF}(A)\wedge\neg\mathrm{IsTrivial}(A)$;
it is not an axiom or a hypothesis of the final classification.

\section{Formal verification and trust boundary}
\label{sec:trust}

The project uses the following pinned Lean and mathlib sources
\cite{lean-source,mathlib-source}. The displayed identifiers also appear
in the project metadata and Lake dependency manifest.
\begin{quote}
\small
\begin{verbatim}
Lean version: 4.34.1
Lean commit:
5045d0056413266e57c625dcd7c365b10e377c52
mathlib commit:
d13f23b723b8a846827a245b89c10fc7d3f11612
\end{verbatim}
\end{quote}

The completed statements are
\path{GPT6Astra.every_order_ten_shf_is_trivial} and
\path{GPT6Astra.not_openProblem10} in \path{MainTheorem.lean}.
The structural, normalization, finite-orbit, and semantic bridge proofs
use the standard foundational axioms \path{propext},
\path{Classical.choice}, and \path{Quot.sound}. The final classification
additionally depends on the evaluation axioms generated by the standard
\path{bv_check} calls for the four finite cases. The actual dependency
lists are printed by \path{#print axioms} in the case modules and in
\path{MainTheorem.lean}, and retained in the verification logs.
No proof is admitted using \path{sorry}, and no user-written axiom
asserts a combinatorial or computational conclusion.

The additional evaluation dependency is substantive. Lean's
normalization and reflection code connects the matrix expressions to
Boolean and bit-vector expressions and then to the certificate problem
\cite{lean-source}. Its standard library proves soundness of the LRAT
checker \cite{lrat-source}. In the standard \path{bv_check} path used
here, the concrete checker calculation is performed by compiled
evaluation and is imported into the proof through generated evaluation
axioms. Consequently, reproduction trusts that evaluation path,
including the compiler and runtime that execute the checker, in
addition to the usual Lean foundations and kernel. The proof is not
claimed to reconstruct every certificate step solely through kernel
reduction. The external SAT search program is not trusted to decide
truth: a saved certificate must pass the checker, and replay does not
require another solver search.

For corroboration, a separately generated Python CNF encoding of each
of the four anchored cases was also refuted by certificate-derived
Lean proof terms whose reported axiom dependency was only
\path{propext}. Those independent audits establish unsatisfiability of
the precise CNFs they inspect. They supplement the final matrix proof;
they neither substitute for its semantic-to-Boolean-to-certificate
chain nor remove its stated compiled-evaluation dependency.

\subsection{Reproduction and immutable proof data}
The stable replay files are \path{certificates/caseK.lrat} for
$K=2,3,4,5$, relative to the Lean project directory. The distribution's
\path{SHA256SUMS} records the actual hashes of the final sources and
certificates. Reproduction should use those recorded source and certificate bytes. The checked source modules use \path{bv_check} with
these stable paths. The generator in
\path{research/direct-matrix/generate_compact_cases.py} records the
indexing scheme and can reproduce the source layout, but it does not
supply an axiom or bypass Lean verification.

With the pinned toolchain and dependencies installed, enter the
\path{gpt6_astra} project directory. Its \path{./lake} wrapper selects
the project toolchain. The following commands replay the cases
sequentially, limiting peak memory use, and then build the integrated
theorem:
\begin{verbatim}
sha256sum -c SHA256SUMS
./lake build FiniteCase2
./lake build FiniteCase3
./lake build FiniteCase4
./lake build FiniteCase5
./lake build MainTheorem
\end{verbatim}
Running \path{./lake build MainTheorem} after these case builds reuses
the checked dependencies. A successful replay means the commands
terminate successfully and the final printed dependencies contain no
\path{sorryAx}; the generated compiled-check evaluation axioms are
expected and explicitly disclosed above. The proof certificate logs
and hash manifest identify the exact release being reproduced.

\appendix
\section{Correspondence with the Lean development}
\label{app:correspondence}

All listed declarations are in namespace \path{GPT6Astra}. The table
maps the mathematical proof and computational interface to the checked
source files.

\small
\begin{longtable}{@{}>{\raggedright\arraybackslash}p{0.29\textwidth}>{\raggedright\arraybackslash}p{0.65\textwidth}@{}}
\textbf{Mathematical statement} & \textbf{File and declarations}\\
\hline
\endfirsthead
\textbf{Mathematical statement} & \textbf{File and declarations}\\
\hline
\endhead
Definitions~\ref{def:separation} and \ref{def:trivial}
& \path{CryptographyResearch.lean}: \path{BinaryMatrix},
\path{Separates}, \path{IsSHF}, \path{IsTrivial},
\path{OpenProblem10}.\\[0.7em]
Row-complement invariance
& \path{CryptographyResearch.lean}:
\path{separates_complementRows_iff},
\path{isSHF_complementRows_iff}.\\[0.7em]
Lemmas~\ref{lem:private-unique} and \ref{lem:elementary-separator}
& \path{MatrixStructure.lean}: \path{PrivateAt},
\path{privateAt_unique},
\path{privateAt_of_elementary_separatesPair}.\\[0.7em]
Theorem~\ref{thm:elementary} and Corollary~\ref{cor:nontrivial-row}
& \path{MatrixStructure.lean}:
\path{isTrivial_of_all_rows_elementary},
\path{exists_non_elementary_row}.\\[0.7em]
Oriented row sets and a nontrivial row
& \path{CutFamilies.lean}: \path{normalizeCut},
\path{normalizeCut_canonical},
\path{selectedCuts_has_minimal_nontrivial}.
The selected cuts are deduplicated sets used to choose a witness;
this does not assume distinct rows in the original matrix.\\[0.7em]
Lemma~\ref{lem:first-row}
& \path{MatrixNormalization.lean}: \path{normalizedMatrix_isSHF},
\path{permuteMatrix_isSHF}, \path{exists_perm_to_initialCut},
\path{exists_normalized_matrix}.\\[0.7em]
Lemma~\ref{lem:anchor}
& \path{MatrixNormalization.lean}: \path{AnchorZero},
\path{anchoredMatrix_isSHF}, \path{anchoredMatrix_fixedRow},
\path{anchoredMatrix_anchorZero},
\path{exists_anchor_normalized_matrix}.\\[0.7em]
Lemma~\ref{lem:double-lex}
& \path{DoubleLex.lean}: \path{matrixLexKey},
\path{rowSwap_key_lt}, \path{columnSwap_key_lt},
\path{exists_doubleLex_anchored},
\path{exists_doubleLex_anchored_fixed}.\\[0.7em]
Theorem~\ref{thm:canonical}
& \path{CanonicalMatrix.lean}:
\path{exists_anchored_canonical_matrix}. The conclusion omits
\path{Shape}; see Remark~\ref{rem:no-weight}.\\[0.7em]
The 840 separation conditions
& \path{FiniteCaseSupport.lean}: \path{separatesB},
\path{separationClause}, \path{separationClause_of_isSHF}.\\[0.7em]
The Boolean lexicographic conditions
& \path{LexBoolean.lean}: \path{prefixB_eq_true_iff},
\path{rowLexClause_of_doubleLexRows},
\path{columnPrefixB_eq_true_iff},
\path{columnLexClause_of_doubleLexAnchoredColumns}.\\[0.7em]
The four finite-case refutations
& \path{FiniteCase2.lean} through
\path{FiniteCase5.lean}; statements
\path{anchored_case_2} through \path{anchored_case_5}.\\[0.7em]
The unconditional classification
& \path{MainTheorem.lean};
\path{every_order_ten_shf_is_trivial} and
\path{not_openProblem10}.\\
\hline
\end{longtable}
\normalsize

The paper also proves the converse direction in
Proposition~\ref{prop:exact} to make the meaning of the formulas
transparent. The formal classification only needs its forward
direction: a nontrivial separating matrix would satisfy the finite
constraints. Auxiliary pair-separation lemmas and alternative cut-family
SAT encodings are not assumptions of the anchored four-case argument
presented here.

\section{Exact computational proof method}
\label{app:computation}

The formal computation operates on the matrix expressions directly;
the separately generated Python CNFs are not inputs to the final
classification proof. The following steps describe the complete
computational part of each case module.

\paragraph{Index generation.}
For fixed $k$, the source generator emits the first-row conditions in
column order $0,\ldots,9$ and the anchor conditions in row order
$1,\ldots,9$. It next emits row-prefix conditions for $i=1,\ldots,8$
and $c=0,\ldots,9$. Column-prefix conditions are generated from
adjacent pairs within the actual lists $[1,\ldots,k-1]$ and
$[k,\ldots,9]$, with prefix rows $1,\ldots,9$. Finally it enumerates
each designated column $a=0,\ldots,9$ and each increasing triple from
$\I\setminus\{a\}$. Thus all case-specific indices are recomputed;
they are not obtained by replacing numerals inside another case's
formula. The same order is retained when replaying a certificate.

\paragraph{Semantic proof blocks.}
The 1002 conditions are grouped into 31 blocks of 32 conditions and a
last block of ten. Each block is a balanced Boolean conjunction.
Every leaf has a proof obtained from a fixed-entry assumption or from
one of the semantic lemmas listed in Appendix~\ref{app:correspondence}.
The local helper \path{bool_and_intro} combines proofs of
$a=\mathrm{true}$ and $b=\mathrm{true}$ into a proof of
$(a\mathbin{\&\&}b)=\mathrm{true}$. The 32 block proofs are then
combined by another balanced conjunction. This layout avoids a local
context with more than a thousand independent hypotheses without
changing the conjunction being proved.

\paragraph{Expansion and canonical atoms.}
The definitions of the combined constraint, its blocks, and the leaf
Boolean operations are unfolded. The tactic
\texttt{canonical\_matrix\_atoms A at hc} visits the resulting expression.
For a concrete matrix application $A\,r\,c$, it computes the values of
its two finite indices and replaces the application with the cached
canonical expression for that pair. It also reduces closed Boolean
\path{decide} expressions, including the numerical comparisons in
fixed-row and prefix conditions. This step is needed because distinct
syntactic representations of the same finite index can otherwise be
abstracted as different Boolean atoms. It does not add an equality as
an assumption: Lean's \path{changeLocalDecl} is called with its
definitional-equality check enabled. Acceptance therefore confirms that
the rewritten hypothesis has the same meaning as the original.

\paragraph{Certificate replay.}
The replay call uses the standard Boolean decision procedure with the
following configuration, where $K$ is the fixed case number:
\begin{verbatim}
bv_check
  (config := { timeout := 120, maxSteps := 2000000,
    structures := false, enums := false, fixedInt := false,
    embeddedConstraintSubst := false })
  "certificates/caseK.lrat"
\end{verbatim}
The generated source writes the configuration on one line and escapes
the field name \path{enums} as required by its Lean syntax; the
multiline display above is explanatory. The configuration retains the
same normalization used for certificate generation, while disabling
passes for structures, enumerations, fixed-width machine integers,
and embedded-constraint substitution that are unnecessary for this
explicit Boolean problem. The step budget controls simplification;
it is not a mathematical restriction on possible matrices. The
\path{timeout} field is inherited from the search configuration and
does not turn an incomplete check into a proof.

The resulting theorem \path{anchored_case_K} concludes false from the
five semantic properties of its input matrix. Combining the four
checked case theorems with
\path{exists_anchored_canonical_matrix}, and splitting the integer
interval $2\leq k\leq5$, proves
\path{every_order_ten_shf_is_trivial}. Unpacking the existential
\path{OpenProblem10} and applying that theorem proves
\path{not_openProblem10}. There is no additional hypothesis concerning
the solver, the certificate, regularity, or symmetry in either final
mathematical statement.

\begin{thebibliography}{9}

\bibitem{designs2015}
C.~Guo, D.~R.~Stinson, and T.~van Trung,
``On Symmetric Designs and Binary 3-Frameproof Codes,''
in C.~J.~Colbourn, ed., \emph{Algebraic Design Theory and Hadamard
Matrices}, Springer Proceedings in Mathematics \& Statistics 133,
Springer, 2015, pp.~125--136.
\href{https://doi.org/10.1007/978-3-319-17729-8_10}{doi:10.1007/978-3-319-17729-8\_10}.
Author manuscript:
\url{https://cs.uwaterloo.ca/~dstinson/papers/sdesign-5.pdf}.

\bibitem{bounds2015}
C.~Guo, D.~R.~Stinson, and T.~van Trung,
``On tight bounds for binary frameproof codes,''
\emph{Designs, Codes and Cryptography} \textbf{77}(2--3) (2015),
301--319.
\href{https://doi.org/10.1007/s10623-015-0037-y}{doi:10.1007/s10623-015-0037-y}.
Preprint: \url{https://arxiv.org/abs/1406.6920}.

\bibitem{wide2024}
Y.~Zhao and X.~Zhang,
``Improved Upper Bounds for Wide-Sense Frameproof Codes,''
\emph{IEEE Transactions on Information Theory} \textbf{70}(12)
(2024), 8636--8646.
\href{https://doi.org/10.1109/TIT.2024.3411150}{doi:10.1109/TIT.2024.3411150}.
Version inspected in the literature audit:
\url{https://arxiv.org/html/2402.05596v1}.

\bibitem{sunwang2025}
C.~Sun and X.~Wang,
``New upper bounds for wide-sense frameproof codes,''
\emph{Designs, Codes and Cryptography} \textbf{93}(8) (2025),
3069--3082.
\href{https://doi.org/10.1007/s10623-025-01631-x}{doi:10.1007/s10623-025-01631-x}.
Only the abstract and first-page preview were available in the
literature audit reported here.

\bibitem{lean-source}
Lean FRO and the Lean contributors,
\emph{Lean 4}, version 4.34.1, official source distribution.
Boolean decision-procedure frontend:
\url{https://github.com/leanprover/lean4/blob/v4.34.1/src/Lean/Meta/Tactic/BVDecide/Main.lean}.

\bibitem{mathlib-source}
The mathlib community,
\emph{mathlib4}, version 4.34.1, official release and source repository.
\url{https://github.com/leanprover-community/mathlib4/releases/tag/v4.34.1}.
The exact dependency commit used here is recorded in
Section~\ref{sec:trust}.

\bibitem{lrat-source}
J.~Clune and the Lean contributors,
\emph{LRAT checker soundness}, Lean standard-library source,
version 4.34.1.
\url{https://github.com/leanprover/lean4/blob/v4.34.1/src/Std/Tactic/BVDecide/LRAT/Internal/LRATCheckerSound.lean}.

\end{thebibliography}
\end{document}
